\section{References}

\begingroup
\footnotesize
\linespread{0.98}\selectfont
\setlength{\parskip}{0pt}
\setlist[enumerate]{itemsep=0pt, topsep=0pt, parsep=0pt, partopsep=0pt}

\subsection{Foundational and early documentary witnesses}

\begin{enumerate}
  \item Miguel Sánchez, \emph{Imagen de la Virgen María Madre de Dios de Guadalupe, milagrosamente aparecida en la ciudad de México} (Mexico City, 1648). First located full printed Spanish narrative and a primary witness to its Apocalypse-centered theological reception.

  \item Luis Laso de la Vega, \emph{Huei tlamahuiçoltica} (Mexico City, 1649), especially the \emph{Nican mopohua}. \href{https://www.loc.gov/item/2021666123/}{Public-domain facsimile and catalog record, Library of Congress}. For navigation and content collation, the study uses the Basilica of Guadalupe's \href{https://virgendeguadalupe.org.mx/el-relato/}{numbered Nahuatl and Spanish presentation, paragraphs 1--218}; the page credits the Spanish translation to Msgr José Luis Guerrero Rosado. The English prose here is an original analytical paraphrase, not a reproduction of that modern translation or a claimed new translation from Classical Nahuatl.

  \item Alonso de Montúfar, \emph{Información} concerning Francisco de Bustamante's sermon of 8 September 1556, printed edition (Madrid: Imprenta de La Guirnalda, 1888). \href{https://bibliotecamexico-fondoreservado.cultura.gob.mx/fondo-reservado/informacion-que-el-arzobispo-de-mexico-d-fray-alonso-de-montufar-mando-a-practicar-con-motivo-de-un-sermon-que-en-la-fiesta-de-la-natividad-de-nuestra-senora-8-de-septiembre-de-1556-predico-en-la-c/}{Catalog and digital-object record, Fondo Reservado de la Biblioteca de México}.

  \item Fortino Hipólito Vera, ed., \emph{Informaciones sobre la milagrosa aparición de la Santísima Virgen de Guadalupe, recibidas en 1666 y 1723} (Amecameca: Imprenta Católica, 1889): \href{https://catalog.hathitrust.org/Record/008398996}{HathiTrust full-view catalog record}; Eduardo Chávez Sánchez, ed., \emph{La Virgen de Guadalupe y Juan Diego en las informaciones jurídicas de 1666: con facsímil del original}, 2nd ed. (2002): \href{https://books.google.com/books/about/La_Virgen_de_Guadalupe_y_Juan_Diego_en_l.html?id=KZRBHQAACAAJ}{bibliographic record}. The complete modern facsimile was not collated; the inquiry's structure and witness problems are controlled by Matovina's critical analysis below.

  \item Bernardino de Sahagún, “Sobre supersticiones,” appendix to Book XI of the \emph{Historia general de las cosas de Nueva España}: \href{https://bibliotecadigital.ilce.edu.mx/sites/fondo2000/vol2/02/htm/sec_5.htm}{official ILCE transcription}. Its Tepeacac/Tepeaquilla, \emph{Tonantzin}, pilgrimage, and Guadalupe notice is used as a polemical primary colonial witness, not neutral ethnography.
\end{enumerate}

\subsection{Papal, liturgical, and official ecclesial sources}

\begin{enumerate}
  \item Second Council of Nicaea as received in the \emph{Catechism of the Catholic Church}, nos.~1159--1162 and 2131--2132, on sacred images; Vatican II, \emph{Dei Verbum}, nos.~2--4, and the \emph{Catechism}, nos.~65--67, on the fullness of Revelation in Christ. \href{https://www.vatican.va/content/catechism/en/part_two/section_one/chapter_two/article_1/iv_sacred_images.html}{Catechism: sacred images}; \href{https://www.vatican.va/archive/hist_councils/ii_vatican_council/documents/vat-ii_const_19651118_dei-verbum_en.html}{\emph{Dei Verbum}}. Exact USCCB Second Edition passage records were separately verified for nos.~1700--1706 on dignity and freedom and no.~2111 on superstition and merely external performance.

  \item Vatican II, \href{https://www.vatican.va/archive/hist_councils/ii_vatican_council/documents/vat-ii_const_19641121_lumen-gentium_en.html}{\emph{Lumen gentium}}, nos.~56 and 60--62; \emph{Catechism of the Catholic Church}, no.~970; and DDF, \href{https://www.vatican.va/roman_curia/congregations/cfaith/documents/rc_ddf_doc_20251104_mater-populi-fidelis_en.html}{\emph{Mater Populi Fidelis}} (4 November 2025), especially nos.~24--28 and 35--38, on Christ's unique mediation and Mary's wholly received, subordinate, maternal cooperation and intercession.

  \item Pius X, apostolic letter \emph{Gratiae quae} (16 June 1910), recording the history of Mexican patronage and Benedict XIV's confirmation: \href{https://www.vatican.va/content/pius-x/la/apost_letters/documents/hf_p-x_apl_19100616_gratiae-quae.html}{official Latin text}.

  \item Pius XII, \href{https://www.vatican.va/content/pius-xii/es/speeches/1945/documents/hf_p-xii_spe_19451012_guadalupe-mexico.html}{radio message for the fiftieth anniversary of the canonical coronation} (12 October 1945).

  \item John Paul II, \href{https://www.vatican.va/content/john-paul-ii/en/homilies/1979/documents/hf_jp-ii_hom_19790127_messico-guadalupe.html}{homily at the Basilica of Guadalupe} (27 January 1979), placing his first international apostolic journey and the Latin American episcopal conference under Mary's protection.

  \item John Paul II, \href{https://www.vatican.va/content/john-paul-ii/es/homilies/1990/documents/hf_jp-ii_hom_19900506_citta-del-messico.html}{homily for the beatification of Juan Diego} (6 May 1990), an official Spanish witness to the beatification and its pastoral interpretation.

  \item John Paul II, post-synodal apostolic exhortation \href{https://www.vatican.va/content/john-paul-ii/en/apost_exhortations/documents/hf_jp-ii_exh_22011999_ecclesia-in-america.html}{\emph{Ecclesia in America}}, no.~11 (22 January 1999).

  \item Congregation for Divine Worship and the Discipline of the Sacraments, \href{https://www.vatican.va/roman_curia/congregations/ccdds/documents/rc_con_ccdds_doc_20000628_guadalupe_sp.html}{decree concerning the feast of Our Lady of Guadalupe throughout America} (25 March 1999; Holy See presentation dated 28 June 2000).

  \item John Paul II, \href{https://www.vatican.va/content/john-paul-ii/en/homilies/2002/documents/hf_jp-ii_hom_20020731_canonization-mexico.html}{homily for the canonization of Juan Diego Cuauhtlatoatzin} (31 July 2002).

  \item Benedict XVI, \href{https://www.vatican.va/content/benedict-xvi/en/homilies/2011/documents/hf_ben-xvi_hom_20111212_america-latina.html}{homily for Latin America on the solemnity of Our Lady of Guadalupe} (12 December 2011).

  \item Francis, \href{https://www.vatican.va/content/francesco/en/homilies/2016/documents/papa-francesco_20160213_omelia-messico-guadalupe.html}{homily at the Basilica of Guadalupe} (13 February 2016); \href{https://press.vatican.va/content/salastampa/en/bollettino/pubblico/2023/12/12/231212d.html}{Guadalupan celebration} (12 December 2023); \href{https://www.vatican.va/content/francesco/es/homilies/2024/documents/20241212-omelia-guadalupe.html}{homily} (12 December 2024).

  \item Leo XIV, \href{https://www.vatican.va/content/leo-xiv/en/homilies/2025/documents/20251212-messa-guadalupe.html}{homily for the feast of Our Lady of Guadalupe} (12 December 2025).

  \item Dicastery for the Doctrine of the Faith, \href{https://www.vatican.va/roman_curia/congregations/cfaith/documents/rc_ddf_doc_20240517_norme-fenomeni-soprannaturali_en.html}{\emph{Norms for Proceeding in the Discernment of Alleged Supernatural Phenomena}} (17 May 2024; effective 19 May). Used only to state current procedure and prevent retroactive terminology.

  \item Insigne y Nacional Basílica de Santa María de Guadalupe, \href{https://virgendeguadalupe.org.mx/documentos-historicos/}{documentary-history portal}. Used as the shrine's own presentation; disputed attributions are cross-checked rather than silently adopted.

  \item Primate Archdiocese of Mexico, \href{https://arquidiocesismexico.org.mx/basilica-de-guadalupe/}{“Basílica de Guadalupe”}, used for current archdiocesan jurisdiction, image custody, and living pilgrimage status. It is not treated as a supernatural-origin decree.
\end{enumerate}

\subsection{Historical, literary, Indigenous, and art-historical studies}

\begin{enumerate}
  \item Ryan Crewe, \href{https://www.cambridge.org/core/journals/americas/article/building-in-the-shadow-of-death-monastery-construction-and-the-politics-of-community-reconstitution-in-sixteenthcentury-mexico/D80E0FD9A9F7C3F90670A6025283C339}{“Building in the Shadow of Death: Monastery Construction and the Politics of Community Reconstitution in Sixteenth-Century Mexico,”} \emph{The Americas} 75, no.~3 (2018), 489--523, DOI 10.1017/tam.2018.31; Cheryl English Martin, \href{https://novohispana.historicas.unam.mx/index.php/ehn/article/download/3356/2911/3307}{“Modes of Production in Colonial Mexico,”} especially 108--109. Used for demographic catastrophe, tribute and labor burdens, contested mission construction, litigation, negotiation, and Indigenous appropriation of institutions.

  \item Getty Research Institute, \href{https://florentinecodex.getty.edu/es/book/12/folio/53r}{\emph{Digital Florentine Codex}, Book 12, fol.~53r}, in the collaborative digital edition of Sahagún's manuscript, used as a primary Nahua-Spanish witness to epidemic memory and as evidence of Indigenous authorship and visual-historical agency.

  \item Timothy Matovina, \href{https://culturasvisuales.org/wp-content/uploads/2017/02/the-origins-of-the-guadalupe-tradition-in-mexico.pdf}{“The Origins of the Guadalupe Tradition in Mexico,”} \emph{Catholic Historical Review} (2013), especially 255--259. Used for the 1563 Verdugo will, Juan Bautista's 1566 procession entry, and a critical reading of the twenty witnesses and leading interrogatory in 1665--1666. His argument is a scholarly interpretation, not a competent ecclesial act.

  \item Rodrigo Martínez Baracs, \href{https://historicas.unam.mx/publicaciones/publicadigital/libros/homenaje/04_14_notas_elaboracion.pdf}{“Notas sobre la elaboración del \emph{Nican mopohua},”} in a UNAM historical volume (2016). Especially important for mapping the authorship and dating disagreement without claiming a false consensus.

  \item Austin Cruz, \href{https://www.cambridge.org/core/journals/horizons/article/huei-tlamahuicoltica-the-hagiographic-origin-of-the-guadalupan-narrative-and-spanish-ideals-for-nahua-sanctity/07CE77AEEDD8A1FB69EB86B5BD5264F1}{“\emph{Huei tlamahuiçoltica}: The Hagiographic Origin of the Guadalupan Narrative and Spanish Ideals for Nahua Sanctity,”} \emph{Horizons} (2025). Open-access hagiographic and audience analysis; one scholarly interpretation, not an ecclesial decree.

  \item Jeanette Favrot Peterson, \href{https://www.cambridge.org/core/journals/americas/article/abs/creating-the-virgin-of-guadalupe-the-cloth-the-artist-and-sources-in-sixteenthcentury-new-spain/71AAA4630ADE9C2A43888FB46982FF31}{“Creating the Virgin of Guadalupe: The Cloth, the Artist, and Sources in Sixteenth-Century New Spain,”} \emph{The Americas} 61, no.~4 (2005). Used for art-historical questions, reported interventions, and human-craft hypotheses; not treated as a technical examination of every layer.

  \item Gisela von Wobeser, \href{https://ru.historicas.unam.mx/handle/20.500.12525/434?show=full}{“Mitos y realidades sobre el origen del culto a la Virgen de Guadalupe,”} especially 150--151, and \href{https://unamglobal.unam.mx/2020/12/11/}{UNAM's public account of the 1556 controversy}. Used as critical historical synthesis of Tepeyac's layered sacred landscape and the early cult, not as a primary transcript.

  \item Lidia E. Gómez García and Eduardo Ángel Cruz, \href{https://www.scielo.org.mx/pdf/ehn/n61/0185-2523-ehn-61-3.pdf}{“El discurso de la desunión: la disputa jurisdiccional por las limosnas de la virgen de Guadalupe en Nueva España, 1572--1607,”} \emph{Estudios de Historia Novohispana} 61 (2019), 3--48, DOI 10.22201/iih.24486922e.2019.61.63139. Used for the \emph{neci} notices, early-cult institutional conflict, alms, work, and jurisdiction.

  \item Guillermo Hurtado Pérez, \href{https://novohispana.historicas.unam.mx/index.php/ehn/article/view/63244}{“La idea de la historia en \emph{Imagen de la Virgen María} de Miguel Sánchez,”} \emph{Estudios de Historia Novohispana} 59 (2018), 74--87. Used to distinguish Sánchez's prodigious, prophetic, and patria narratives and to identify his Apocalypse-centered historical method.
\end{enumerate}

\subsection{Technical and claim-audit sources}

\begin{enumerate}
  \item Arturo Rocha Cortés, \href{https://revistas.uic.mx/index.php/voces/es/issue/download/134/95}{“Incidente del ácido derramado ca.~1784 sobre el Sagrado Original de Nuestra Señora de Guadalupe,”} \emph{Voces} 37 (2012), presenting and analyzing AHBG, Correspondencia con el Supremo Gobierno, caja 3, exp.~54. Used for the later 1820--1823 inquiry and its uncertainties; the publication does not repeat “1791 nitric-acid accident” as a secure fact.

  \item \emph{El Universal}, \href{https://www.eluniversal.com.mx/cultura/el-atentado-hacia-la-virgen-de-guadalupe-de-hace-100-anos/}{“El atentado hacia la Virgen de Guadalupe que desató la ira de fieles hace 100 años”} (25 December 2021), republishing its reporters' 15 November 1921 account. Used only for the observed altar, candlestick, crucifix, frame, image, and glazing conditions after the explosion; not for a technical impossibility or culprit finding.

  \item Philip S. Callahan, \emph{The Tilma under Infra-red Radiation: An Infrared and Artistic Analysis of the Image of the Virgin Mary in the Basilica of Guadalupe} (Washington, DC: Center for Applied Research in the Apostolate, 1981). \href{https://search.worldcat.org/title/11200397}{WorldCat bibliographic record}. Used with its limited scope; popular paraphrases are not attributed to the report without verification.

  \item NASA public-site negative search for \emph{Guadalupe}, \emph{tilma}, and claimed agency examination, repeated 16 July 2026. No relevant agency report was located. A dated negative search supports rejection of the attribution, not a universal claim that no individual with an aerospace connection ever commented.

  \item The material object, its dated reproductions, and published conservation and art-historical descriptions. No direct access, sampling, imaging, or new experiment was performed for this monograph. Technical claims therefore remain bounded by the methods and access of their actual reports.
\end{enumerate}

\endgroup
